
\paragraph{Aggregated Secrecy.}
Unlike in the plain VSS setting,
an \AVSS allows for aggregating shares and commitments that could both be honestly as well as maliciously generated.
Hence, we require that the \AVSS fulfills the notion of \emph{aggregated secrecy},
which we formalize in Figure~\ref{fg:avss-secrecy}.


\begin{remark}[Aggregated secrecy for an \AVSS is a generalization of VSS secrecy]
  The secrecy game for a VSS scheme as shown in Figure~\ref{fg:vss-secrecy} only allows the adversary to receive a single set of shares and their corresponding commitment from the environment.
  Hence, the aggregated secrecy game in Figure~\ref{fg:avss-secrecy} is a strict generalization of the VSS secrecy game,
  as the adversary in the secure aggregation game can not only query for many shares and commitments, but submit its own.
  As a consequence, any \AVSS that fulfills aggregated secrecy is also secret.
\end{remark}


\begin{figure}[p]
  \begin{pchstack}[boxed,center,space=1em]
  \begin{pcvstack}
    \procedure[]{$\AggSecrecyExp_{\avss}(\lambda, \secret)$}{
    b \rgets \zo \\
    \text{issued} \gets 0 \\
		\pp \leftarrow  \avss.\Setup(1^\secp) \\
    \vsschal \rgets \avss.\vfunc(\secret) \\
    \issuedchal \gets 0 \\
    \auxstring \rgets \avss.\auxdistr \\
    %\text{\algcom Pick random auxiliary string} \\
    \simtotal \gets 0  \\
    Q_1 \gets \emptyset, Q_2 \gets \emptyset, Q_3 \gets \emptyset  \\
    (n, t, \corrupt, \advstate) \rgets \adv(1^\lambda) \\
    % We already said that n,t are natural numbers so we don't need to check their types.
    \pcif t > n \vee \corrupt \not\subseteq \set{n}  \\
    \textbf{or }\lvert \corrupt \rvert \geq t \\
    \pcind \randcoin \rgets \zo;\ \pcreturn \randcoin \\
    \honest \gets \set{n} \setminus \corrupt \\
    %
    \advstate' \rgets \adv^{\getshare,\recvshare, \Ogettweak}(\advstate) \\
    \pcif \issuedchal = 0 \\
    \pcind \randcoin \rgets \zo;\ \pcreturn \randcoin \\
    \pcif b=1 \text{ and } \text{issued} = 0 \\
    %\pcind (\secretone^*,\vsscommit^*) \gets Q_1 \\
    %\pcind \{(\{(i,\share_{ji})\}_{i \in \honest},\vsscommit_j) \}_{j \in \set{\simtotal}} \gets Q_2 \\
    \pcind \mathsf{in} \gets ( \{\vsscommit_j \}_{j \in \set{\simtotal}} \cup \{\vsscommit^* \}, \auxstring) \\
    \pcind \avsscommit \gets \avss.\aggregatepub(\mathsf{in}) \\
    \pcind \pcif \avss.\derivepshare(0, \avsscommit) \neq \vsschal \\
    \pcind \pcind \pcreturn 1 \\
    b' \rgets \adv^{\Ogettweak}(\advstate', \auxstring) \\
    \text{\algcom $\adv$ performs} \\
    \text{\algcom  aggregation directly} \\
    \pcreturn 1\ \pcif b' = b \\
    \pcreturn 0
  }
  \end{pcvstack}
  \pchspace
  \begin{pcvstack}
    \procedure[]{$\getshare() \text{\algcom Performs secret sharing}$}{
      \pcif \issuedchal = 0 \\
      \pcind \issuedchal \gets 1 \\
      \pcind \pcif b=0 \\
      \pcind \pcind \{(j,\share_{j})\}_{j =1}^n, \vsscommit \rgets \avss.\issueshares(s, n, t) \\
      \pcind \pcind Q_2 \gets  Q_2 \cup   \{ (\{(j,\share_{j})\}_{j \in \honest},, \vsscommit) \} \\
      \pcind \pcind \pcreturn ( \{(j,\share_j)\}_{j \in \corrupt}, \vsscommit) \\
      \pcind \pcelse \text{\algcom $b=1$ case} \\
      \pcind \pcind \mathbf{out} \rgets \simshare(\lambda, \corrupt, \vsschal)\\
      \pcind \pcind (\secretone^*, \{(j,\share_j) \}_{j \in \corrupt},
      \vsscommit^*) \gets \mathbf{out} \\
      \pcind \pcind Q_1 \gets (\secretone^*, \vsscommit^*) \\
      \pcind \pcreturn ( \{(j,\share_j)\}_{j \in \corrupt}, \vsscommit^*) \\
      \simtotal \gets \simtotal + 1;\
      s_\simtotal \rgets \secretspace \\
      \{(j,\share_{\simtotal j})\}_{j =1}^n, \vsscommit_\simtotal \rgets \avss.\issueshares(s_\simtotal, n, t) \\
      Q_2 \gets  Q_2 \cup   \{ (\{(j,\share_{\simtotal j})\}_{j \in \honest},, \vsscommit_\simtotal) \} \\
      \pcreturn (\{(j,\share_{\simtotal j})\}_{j \in \corrupt}, \vsscommit_\simtotal)
  }
  \pcvspace
    \procedure[]{$\recvshare(\{ (j, \share_{\simtotal j}) \}_{j \in \honest}, \vsscommit_\simtotal)$}{
    %\text{\algcom Receives shares from $\adv$} \\
	  \pcfor j \in \honest \\
    \pcind \pcif \avss.\verify(j, \share_{\simtotal j}, \vsscommit_\simtotal) \neq 1 \\
    \pcind \pcind \pcreturn \bot \\
    Q_2 \gets  Q_2 \cup  \{ (\{(j,\share_{\simtotal j})\}_{j \in \honest}, \vsscommit_\simtotal) \} \\
    \simtotal \gets \simtotal + 1 \\
  }
  \pcvspace
  \procedure[]{$\Ogettweak(\pubaggregateset_i, \auxstring_i) $}{
    \pcif b=1 \\
    \pcind \tweak \gets \simgettweak(\pubaggregateset_i, \auxstring_i, Q_1, Q_2) \\
    \pcelse \\
    \pcind \tweak \gets \avss.\gettweak(\pubaggregateset_i, \auxstring_i) \\
    \pcreturn \tweak
}
  \end{pcvstack}
  \end{pchstack}

  \caption{Aggregated secrecy experiment for a zero-indexed Aggregatable Verifiable Secret Sharing scheme (\AVSS).
}
  \label{fg:avss-secrecy}
\end{figure}



When executing the experiment shown in Figure~\ref{fg:avss-secrecy},
%the environment accepts as input the security parameter $\lambda$,
%a secret $\secret$ from the secret space $\secretspace$, and a bit $b$.
%The environment then interacts with the adversary to perform secret sharing and to receive secret shares.
%The adversary is able to receive $\lvert \corrupt \rvert$ number of secret shares for a fresh secret by querying the oracle $\getshare$,
%and send $\lvert \honest \rvert$ number of its own secret shares to the environment by querying $\recvshare$ with its contributions.
the environment simulates one secret sharing operation with respect to a challenge $\vsschal$ when the random bit $b=1$.
The environment then allows the adversary to query for corrupted parties' shares by $\getshare$,
and submit its own shares for honest parties by $\recvshare$.
The adversary is then provided with the auxiliary string $\auxstring$.
The adversary automatically wins if the aggregated commitment $\avsscommit$ does not
fulfill the check
$\avss.\derivepshare(0, \avsscommit) \neq \vsschal$.
Otherwise,
%the adversary learns sufficient information to perform $\avss.\aggregatepub$ and $\avss.\aggregatepriv$ directly.
$\adv$ outputs a guess $b'$ to guess if it is in the real or simulated environment.

The advantage of an adversary $\adv$ against $\avss$ in the aggregated secrecy experiment as defined in Figure~\ref{fg:avss-secrecy} with respect to algorithms $\simshare,\simgettweak$ is
\[
  \advantaggsecrecy{\avss, \adv, \simshare,\simgettweak}(\lambda) = \max_{\secret \in \secretspace}{ \bigabs{\Pr[\AggSecrecyExp_{\avss, \adv,\simshare,\simgettweak}(\lambda,s) = 1 ] - 1/2 } }
 \]

\begin{definition}
\label{defn:aggsecrecy}
An \AVSS scheme $\avss$ achieves \defn{aggregated secrecy} if
there exist efficient algorithms $\simshare, \simgettweak$ such that
for all probabilistic polynomial-time adversaries $\adv$,
the function $\advantaggsecrecy{\avss, \adv,\simshare,\simgettweak}(\lambda)$
is a negligible function of $\lambda$.
\end{definition}

